\documentclass[11pt,a4paper]{article}

% ------------------------------------------------------------------
% House layout, shared with the sibling preprints in docs/research/:
% 11pt a4, 2.5cm margins, Computer Modern, \maketitle front matter,
% abstract + keywords + JEL codes, numbered sections, table of contents.
% ------------------------------------------------------------------
\usepackage{amsmath, amssymb, amsthm}
\usepackage[a4paper, margin=2.5cm]{geometry}
\usepackage{array}

\theoremstyle{definition}
\newtheorem{definition}{Definition}[section]
\newtheorem{remark}{Remark}[section]
\newtheorem{proposition}{Proposition}[section]
\theoremstyle{plain}
\newtheorem{theorem}{Theorem}[section]
\newtheorem{lemma}{Lemma}[section]

\newcommand{\R}{\mathbb{R}}
\newcommand{\E}{\mathbb{E}}
\newcommand{\N}{\mathbb{N}}
\newcommand{\ind}{\mathbf{1}}
\newcommand{\norm}[1]{\left\lVert #1 \right\rVert}
\newcommand{\inner}[2]{\langle #1,\, #2 \rangle}
\newcommand{\corr}{\operatorname{corr}}
\newcommand{\sgn}{\operatorname{sign}}
\newcommand{\code}[1]{\texttt{\detokenize{#1}}}
\DeclareMathOperator*{\argmax}{arg\,max}

% ============================================================
\begin{document}

\title{\textbf{Exact or Nothing}\\[6pt]
\large Gate-Governed Measurement of the Liquidation Map\\
on a Transparent Perpetual-Futures Venue}

\author{Yigit Onat\\
\textit{Independent Researcher}\\
\texttt{yionat@gmail.com}}
\date{August 2026}
\maketitle

% ============================================================
\begin{abstract}
Research on liquidation cascades is usually conducted on a liquidation map that has been
\emph{estimated}---from aggregate open interest, an assumed leverage distribution, and a reconstructed
population of positions. We ask what becomes possible, and what becomes visible, when the map is
instead \emph{observed}: on Hyperliquid, per-position liquidation prices are published, the mark price
is built from an oracle the venue itself publishes, and any address's fill history is public. We
describe an instrument built on two rules---\emph{exact or nothing} (a quantity that cannot be had
exactly is not computed, because the approximation is what makes a backtest lie) and \emph{gates
before models} (every command downstream of a validation gate refuses to run when that gate is
missing, stale, or failing, and the refusal is the product). The instrument's substrate is a
write-once capture store in which every record carries both an event timestamp and an ingest
timestamp, making lookahead a queryable property rather than a code-review hope, and a hash-chained
ledger in which a failed verdict cannot quietly leave the record. We then report what the instrument
measured. Its feed gate passes twenty checks; its map gate does not pass, and we explain why that is
the informative outcome. Registry coverage is $24.7$--$38.0\%$ of venue open interest by symbol under
a conservative side convention; the derived liquidation price reproduces the venue's published value
to within $10^{-4}$ for $40.8\%$ of $2{,}926$ positions, with an unexplained residual we decline to
model away. The central finding is negative and, we believe, general to any venue where positions must
be sampled: of $48$ observed liquidations, the notional-seeded registry contained $2.9\%$ of the
liquidated \emph{wallets} but $21.9\%$ of the liquidated \emph{notional}. Mass and events are carried
by different populations, so a sampling scheme optimized for one is mis-specified for the other, and
no shape-preserving rescaling repairs the difference. We report that the mapped fraction is currently
$0.0\%$, show that this figure is confounded by snapshot ordering rather than informative, and specify
the scheduled experiment that discriminates among the three explanations. We present no performance
figures; the contribution is an instrument, a discipline of refusal, and four measurements that
change what the observational obstacle is understood to be.
\end{abstract}

\smallskip
\noindent\textbf{Keywords:} market microstructure, liquidation cascade, perpetual futures,
decentralized exchange, measurement, research infrastructure, causality, reproducibility, coverage
bias, sampling

\smallskip
\noindent\textbf{JEL Classification:} C55, C81, G14, G17

\tableofcontents
\bigskip

% ============================================================
\section{Introduction}
\label{sec:intro}
% ============================================================

A leveraged perpetual-futures position carries a deterministic liquidation price. In aggregate these
prices form a spatial field---a \emph{liquidation map}---and the structural claim, developed
formally in the companion paper [10], is that when price reaches a dense region of this
field, forced market liquidations impose impact that can carry price to the next dense region,
producing a self-reinforcing cascade. The claim is intuitive, mechanically plausible, and hard to
test, for a reason that is observational rather than theoretical: on most venues the map cannot be
seen. Position-level data is private; what is published is aggregate open interest. The map is
therefore reconstructed by assuming a leverage distribution over that aggregate, and every
downstream result inherits the assumption.

This paper reports on an instrument built to remove the assumption rather than refine it.
Hyperliquid is, to our knowledge, the venue where the fewest such assumptions are necessary: it
publishes a liquidation price per position, it liquidates against a mark price built from an oracle
it also publishes, funding is settled hourly and observable, its backstop vault's inventory is
visible, and the fill history of any address can be read by anyone. On such a venue, the interesting
question is not how to estimate the map better. It is what the map---and the research programme
built on it---looks like when estimation is refused outright.

\subsection{Two rules}

The instrument is organized around two rules, and both are enforced in code rather than intended in
prose.

\paragraph{Exact or nothing.} Where a quantity cannot be obtained exactly, it is not computed. A
position that is not observed is absent from the map and enters the reported coverage deficit; it is
not imputed. There are no estimated maps, no assumed leverage mixes, and no aggregator history. The
justification is asymmetric: an approximation that is 80\% right yields a backtest that is
confidently wrong, whereas an acknowledged hole yields a coverage number, and a coverage number can
gate a decision.

\paragraph{Gates before models.} The research programme is expressed as a small set of experiments
that are each \emph{allowed to end the project}. Each is a command that writes a \textsf{PASS} or
\textsf{FAIL} verdict to an append-only ledger, and every command downstream of a gate refuses to
run while that gate is missing, stale, or failed. The refusal is not a safety feature bolted onto a
research tool; it is the primary output. A system that will happily produce a Sharpe ratio from data
whose integrity check failed has not produced a Sharpe ratio.

\begin{table}[t]
\centering
\small
\begin{tabular}{@{}l l l@{}}
\hline\hline
Gate & Question & What a \textsf{FAIL} kills \\
\hline
\code{feed}        & is the captured data intact and causal?                & everything \\
\code{map}         & does the map predict the venue's own liquidations?     & the map branch \\
\code{magnet}      & does the cluster pull, net of costs, beat a baseline?  & the magnet experts \\
\code{persistence} & does skill in $[t-90d,\,t]$ predict $(t,\,t+30d]$?     & the smart-money branch \\
\code{decay}       & does cohort-flow alpha survive past minutes?           & copyability \\
\hline\hline
\end{tabular}
\caption{The gate ladder. Each gate is permitted to terminate the branch beneath it; two of them are
permitted to terminate the project. Ordering them before the modelling work is deliberate.}
\label{tab:gates}
\end{table}

\subsection{Contributions}

\begin{enumerate}
  \item A capture and validation substrate in which \emph{lookahead is queryable}: every record
        carries the venue's event clock and our ingest clock, and a replay audit fails on any feature
        that read a record it could not have known (Section~\ref{sec:capture}).
  \item A refusal protocol built on a hash-chained ledger, with the property that a failed verdict
        cannot be quietly re-run away, and a formal statement of the condition under which a
        downstream command is permitted to execute (Section~\ref{sec:gates}).
  \item A liquidation map whose coverage is a first-class, printed, gate-thresholded quantity rather
        than an unstated assumption, together with the side-convention question that makes it a
        factor of two (Section~\ref{sec:map}).
  \item Four measurements, of which the important one is negative: the population that holds notional
        and the population that gets liquidated are different, and the gap between them is an order
        of magnitude (Sections~\ref{sec:measurements}--\ref{sec:mismatch}).
\end{enumerate}

We report no trading performance. The instrument's map gate does not currently pass, and under the
second rule that fact bars the downstream work rather than being reported alongside it.

% ============================================================
\section{Why a Single Transparent Venue}
\label{sec:venue}
% ============================================================

Restricting to one venue is usually a limitation. Here it is load-bearing, and the reason is that
the alternative silently reintroduces the estimation the first rule exists to forbid.

A cross-venue design would supply global price pressure from a large centralized exchange and
liquidation geometry from the decentralized one. That design requires a second clock, a second cost
model, and an alignment assumption between them. The single-venue design does not, because the venue
publishes the oracle price against which it constructs its own mark price. The premium
\begin{equation}
  \pi_t \;=\; \frac{p^{\text{mark}}_t - p^{\text{oracle}}_t}{p^{\text{oracle}}_t}
  \label{eq:premium}
\end{equation}
is then a native, exactly observable measure of exactly the global pressure a second-venue sidecar
would have been introduced to supply---and the causal arrow points the right way, since global price
drives liquidations against a map we observe directly rather than infer.

Three further quantities have no clean analogue elsewhere and materially shape what is measurable.
First, funding is settled \emph{hourly}, which at cascade horizons is a first-order carry term rather
than a rounding error. Second, the venue operates a backstop liquidity vault whose inventory change,
$\Delta_t^{\text{HLP}}$, is public; when forced flow overwhelms the resting book the vault absorbs
it, so $\Delta_t^{\text{HLP}}$ \emph{measures} the portion of a cascade the book failed to absorb.
Absorption exhaustion, which elsewhere must be inferred from price behaviour, is here directly
observed. Third, a single address's perpetual, spot, and vault activity are all visible together, so
an intra-venue hedge is detectable rather than an unmodelled confound in any flow-following study.

What is given up, and we state it rather than minimize it: liquidation density on other venues (only
ever an estimate, which the first rule rejects in any case); a long history spanning many regimes,
which is handled with wider embargoes and smaller models and never by borrowing history from another
venue; the thin-alt cascade tail; and the operational risk of depending on one venue, which is
accepted explicitly.

% ============================================================
\section{Capture: Making Lookahead Queryable}
\label{sec:capture}
% ============================================================

Every result downstream depends on the claim that a feature computed for time $T$ used only
information available at $T$. That claim is normally defended by code review. We make it a property
of the data.

\subsection{Two clocks}

Each captured record carries two timestamps: $t^{\text{event}}$, the venue's clock, and
$t^{\text{ingest}}$, the first moment at which we could have known the record. The causality
contract is then a filter rather than a convention.

\begin{definition}[Causal readable set]
\label{def:causal}
For a bar closing at time $T$, the admissible input set is
\[
  \mathcal{F}_T \;=\; \bigl\{\, r \;:\; t^{\text{ingest}}(r) \le T \,\bigr\},
\]
and a feature is \emph{causal} if it is measurable with respect to $\mathcal{F}_T$.
\end{definition}

Definition~\ref{def:causal} is enforced by replay: the audit recomputes features under the filter and
fails on any difference from the production value. The distinction from $t^{\text{event}} \le T$
matters precisely when it is most dangerous---a record the venue stamped before $T$ but which reached
us after $T$ is exactly the kind of input that inflates a backtest while looking defensible.

\subsection{A write-once store}

Raw records are appended to per-stream files that rotate on the UTC hour. When a file closes it is
hashed and an entry is written to a manifest; the file is never reopened. A restart inside the same
hour begins a \emph{new part} rather than appending to the closed one, which gives the store a useful
property: a changed digest always means corruption or tampering and never merely a bounced daemon. A
per-stream sequence number continues across restarts, resumed from the manifest, so a hole in the
sequence means records were \emph{lost} rather than that the process died.

The feed gate audits this structure directly. Its checks include manifest integrity (every closed
file still hashes to its recorded digest), orphan detection (an unmanifested file older than two
hours indicates a writer that died without closing, so its tail may be truncated), per-stream
sequence continuity, ingest monotonicity, gap detection at $60\times$ the stream's cadence hint,
staleness at $20\times$, and a staleness check on the rate-limit constant table itself---each limit
is recorded with the date a human last verified it against the venue's documentation, and the gate
warns once that stamp exceeds ninety days. A limit that moved without anyone noticing is exactly the
failure that stamp exists to surface.

\subsection{Reconstruct, do not poll}

The registry cannot be maintained by polling. Measured on 2026-08-07, a full sweep of clearinghouse
state across $3{,}596$ wallets takes approximately five minutes and consumes the entire per-IP
request-weight budget of $1{,}200$ units per minute, leaving nothing for the trade tape. Polling
frequency would therefore \emph{become} the map's time resolution---an infrastructural constraint
silently determining a scientific one.

Positions are instead maintained between snapshots by reconstruction from the public fill stream,
with polling demoted to reconciliation on a rotating sample, and drift beyond tolerance failing the
feed gate. Cross-margined and isolated-margined positions take separate paths, because a
cross-margined account liquidates as a unit: one position's move relocates that account's entire
contribution to the map. The bootstrap limitation is stated plainly---the tape supplies
\emph{changes} while snapshots supply \emph{levels}, so a wallet is reconstructable only from its
most recent snapshot forward.

% ============================================================
\section{The Gate Protocol and the Refusal Ledger}
\label{sec:gates}
% ============================================================

\subsection{The ledger}

Gate verdicts, specification freezes, run records, and periodic observations are appended to a
hash-chained log. Writing $e_i = (i,\, \tau_i,\, \kappa_i,\, \pi_i)$ for the $i$-th entry's sequence
number, timestamp, kind, and payload, the chain is
\begin{equation}
  h_i \;=\; \mathrm{SHA\text{-}256}\bigl( \mathrm{enc}(i,\, \tau_i,\, \kappa_i,\, \pi_i,\, h_{i-1}) \bigr),
  \qquad h_{-1} = 0^{64},
  \label{eq:chain}
\end{equation}
with $\mathrm{enc}$ a canonical (key-sorted, separator-fixed) JSON encoding so that the digest is
reproducible. Verification walks the chain and reports the first index at which the sequence, the
predecessor link, or the recomputed digest disagrees.

The construction is standard [5,,8] and the threat model is not an adversary.
It is the researcher. Equation~\eqref{eq:chain} makes deletion of an inconvenient result detectable,
so the fourteen variants tried before the one that worked remain on the record. This is the
mechanism by which multiple-testing discipline becomes structural rather than remembered---a concern
developed at length for the feature and process layers in [11,\,13].

\subsection{The refusal predicate}

\begin{definition}[Admissibility]
\label{def:admissible}
Let $V(g)$ denote the most recent verdict for gate $g$, with fields $\mathrm{pass}(g) \in \{0,1\}$
and timestamp $\tau(g)$. A command depending on $g$ may execute at time $t$ if and only if
\[
  V(g) \text{ exists}
  \;\wedge\;
  \mathrm{pass}(g) = 1
  \;\wedge\;
  t - \tau(g) \le \Delta,
\]
with $\Delta$ a freshness limit defaulting to $24$ hours. Otherwise the command raises a refusal
naming which of the three conditions failed.
\end{definition}

The third clause is the one that is easy to omit and expensive to omit. A verdict describes the data
as it stood when the gate ran; past some age it says nothing about the data about to be used. Without
a freshness bound, a gate that passed once becomes a permanent licence.

\begin{remark}[A gate that cannot be evaluated must not pass]
\label{rem:unevaluable}
Some criteria are unanswerable on demand. The map gate's predictive criterion requires map snapshots
\emph{followed by} observed liquidations, so it is unanswerable until capture has accumulated both.
The gate reports insufficient history and \textsf{FAIL}s. Treating an unevaluable criterion as
passing---or as not applicable---is the single most convenient error available in this design, and it
would convert the whole protocol into decoration.
\end{remark}

Remark~\ref{rem:unevaluable} is why the results in Section~\ref{sec:measurements} include a failing
gate and no performance numbers. That is the protocol operating as designed, not a gap in the work.

% ============================================================
\section{The Map and Its Coverage}
\label{sec:map}
% ============================================================

The map bins observed positions by liquidation price relative to the current mark, recording notional
per bucket, the cross-margined share separately, and position counts; band aggregates
$L^{\uparrow}(b)$ and $L^{\downarrow}(b)$ are taken over $b \in \{0.5, 1, 2, 5\}\%$. The geometry and
the cascade mechanics built on it are the subject of [10]; here we are concerned with the
number that qualifies the field.

\begin{definition}[Coverage]
\label{def:coverage}
For coin $c$ at time $t$, with $N_t$ the total notional of observed positions and
$\mathrm{OI}_t$ the venue's reported open interest,
\begin{equation}
  \widehat{\rho}_t \;=\; \frac{N_t}{s \cdot \mathrm{OI}_t},
  \qquad s \in \{1, 2\},
  \label{eq:coverage}
\end{equation}
where $s$ is the side convention.
\end{definition}

The convention deserves its own sentence because it is a factor of two on the one quantity by which
the map is judged. Under the standard one-sided reading, an exchange counts each contract once in
$\mathrm{OI}_t$; but every long has a matching short, and a registry that tracks wallets sees both
sides. Total position notional outstanding is then $2\,\mathrm{OI}_t$, and the conservative choice is
$s=2$. We take it, print it on every map alongside the coverage figure, and list it as an open
verification item rather than resolving it by preference.

A map without its coverage number is, in the design's own phrasing, a lie with a chart on it. The
number is not decoration: the map gate thresholds on it at $\widehat{\rho} \ge 0.25$.

\subsection{Observing realized liquidations}

Scoring the map requires realized liquidation events, and finding them is less obvious than it
appears. The public trade tape carries no liquidation marker. The venue's fill records carry one in
two places: on the liquidated wallet's own fill, in a direction string, and---far more usefully---on
the \emph{counterparty's} fill, as a structured object naming the liquidated user, the mark price,
and the method (market or backstop).

The asymmetry between the two is extreme. Measured on 2026-08-08 over $117{,}485$ fills from $70$
wallets, exactly \emph{one} fill carried the liquidated-side direction string, while $16$ of the $70$
wallets had taken the other side of somebody's liquidation. The efficient observer set is therefore
not the registry at all: it is a handful of high-volume market makers, because whoever absorbs forced
flow sees it. Since the fill endpoint has no per-user cap over the request interface---unlike the
streaming interface's limit of fifteen users per connection---this scales with the weight budget
rather than with connection count.

\begin{remark}[The lookahead trap in scoring]
\label{rem:lookahead}
An event that predates the position snapshot it is scored against is not a prediction; it is a
memory. Scoring drops such events, and reports the number dropped rather than absorbing it silently.
Of the first $48$ events observed, $24$ predated the snapshot and were correctly excluded on these
grounds---a figure worth reporting precisely because a scorer without this check would have counted
them.
\end{remark}

% ============================================================
\section{Measurements}
\label{sec:measurements}
% ============================================================

All figures below are read from the instrument's ledger and are reported with the date on which the
verdict or observation was appended. Nothing in this section is a backtest.

\begin{table}[t]
\centering
\begin{tabular}{@{}l l l@{}}
\hline\hline
Quantity & Value & Provenance \\
\hline
\code{gate feed} verdict     & \textsf{PASS}, $20$ checks, $0$ failed        & ledger, 2026-08-07 \\
\code{gate map} verdict      & \textsf{FAIL} on \code{predictive}            & ledger, 2026-08-08 \\
Coverage, BTC                & $0.315$                                       & ledger, $s=2$ \\
Coverage, ETH                & $0.380$                                       & ledger, $s=2$ \\
Coverage, SOL                & $0.247$                                       & ledger, $s=2$ \\
Derivation exact fraction    & $0.408$ of $2{,}926$ positions                & rel.\ err.\ $<10^{-4}$ \\
Derivation median rel.\ err. & $0.152$ (p90: $0.921$)                        & vs.\ published value \\
Liquidated wallets in registry & $2.9\%$ of $35$ wallets                     & ledger observation \\
Liquidated notional in registry & $21.9\%$ of $\$223{,}491$                  & ledger observation \\
Mapped notional fraction     & $0.0\%$ \emph{(confounded, Section~\ref{sec:mismatch})} & ledger observation \\
\hline\hline
\end{tabular}
\caption{Measurements to date. The map gate fails; under Remark~\ref{rem:unevaluable} this is the
correct verdict rather than a missing result.}
\label{tab:measurements}
\end{table}

\subsection{Coverage}

A leaderboard exposes $41{,}392$ wallets with equity and windowed volume---roughly forty times the
population that trade-stream discovery reaches, which is what makes a seeded registry a different
regime rather than a larger version of the same one. The two natural seeds behave differently, and
the difference turns out to matter more than the levels (probe of 2026-08-07): the top $2{,}000$
wallets \emph{by equity} span $61\%$ of venue open interest but only $28\%$ hold any position at a
given moment, while the top $2{,}000$ \emph{by weekly volume} span $39\%$ with $46\%$ active. Their
union reaches $69\%$. Accordingly the map registry seeds by equity, since it wants whoever holds
size, and the skill cohort seeds by volume, since it wants whoever trades.

The operational registry's realized coverage under Definition~\ref{def:coverage} with $s=2$ is $0.315$
for BTC, $0.380$ for ETH, and $0.247$ for SOL. We note without resolving it that the ratio between
these figures and the $69\%$ scoping probe is close to $s$ itself, which is consistent with the probe
having been computed one-sided; the convention remains an open verification item, and we report both
numbers rather than adopting the more flattering one.

\subsection{Fidelity of the derived liquidation price}

The venue publishes a liquidation price per position, and that published value is the source of
truth. A derived value is needed only for positions whose size changed since the last snapshot and
which therefore have no current published number. Deriving is thus an approximation with a measured
error, not a replacement---and measuring the error is the point.

Diffed against the published field over $2{,}926$ positions: median relative error $0.152$, ninetieth
percentile $0.921$, and agreement to within $10^{-4}$ for $40.8\%$. Restricted to accounts holding no
isolated margin, agreement rises to $58\%$ over $1{,}074$ cross positions. The residual is
unexplained. Size-tiered maintenance margin on large positions is the leading suspect, and the map
gate declines to fail on fidelity alone---the derivation is a maintenance approximation, not the
primary source---but it does fail on a \emph{collapse} below $20\%$, which would indicate that the
venue's margin rules had moved while the fill-maintained map continued blind.

We draw attention to this as a small instance of the first rule paying for itself. A design that
derived liquidation prices unconditionally would have shipped a map that is materially wrong for
roughly three positions in five, with nothing in the system able to say so.

% ============================================================
\section{The Population Mismatch}
\label{sec:mismatch}
% ============================================================

Measuring coverage was expected to close the observational problem posed in [10]. It did
not. It exposed a second problem, different in kind, and the paper's central finding is that the
second one is not repairable by the standard remedy for the first.

\subsection{Two distinct biases}

The usual correction for partial observation assumes the observed field is a uniform thinning of the
true one, so that rescaling by $\mathrm{OI}_t / N_t$ restores magnitude while leaving every ratio
feature untouched. The assumption is that the sampling is unbiased in \emph{shape}. That is what
fails here, and it fails for a structural reason:

\begin{quote}
The registry is seeded by equity, because the map wants whoever holds size. But the wallets that get
liquidated are small and highly leveraged. Sampling proportional to notional is precisely \emph{not}
sampling proportional to liquidation events.
\end{quote}

Of $48$ liquidations observed and appended to the ledger, the registry contained $2.9\%$ of the
liquidated \emph{wallets} ($35$ distinct addresses) and $21.9\%$ of the liquidated \emph{notional}
($\$223{,}491$). A second reading one cycle later gave $2.8\%$ and $21.9\%$ over $49$ events. The
order-of-magnitude gap between the two percentages is the second bias stated quantitatively, and it
is invisible to any coverage figure of the form \eqref{eq:coverage}, which is denominated in notional
and therefore reports the flattering one of the two.

The finding cuts both ways, and we resist collapsing it to either. Cascades are driven by notional,
and catching a fifth of liquidated notional is not nothing; per-position scoring is not obviously
dead. What the field is poor at is seeing the long tail of small forced exits---which matters exactly
to the extent that those small exits are what \emph{initiates} a chain rather than what continues one.

\begin{remark}[Coverage of what?]
\label{rem:coverage_of_what}
The generalizable form of this result is that ``insufficient coverage'' is an incomplete diagnosis.
Coverage is a ratio, and a field can be simultaneously adequate for the mass that sustains a cascade
and inadequate for the events that start one. Any venue on which positions must be sampled rather
than enumerated is exposed to this, and a single coverage scalar will not reveal it. We now report
coverage both by notional and by event count for that reason.
\end{remark}

\subsection{Why the reading cannot yet be believed}

Both readings show $0.0\%$ of liquidated notional \emph{mapped}---attributable to a position the map
had marked in advance. Taken at face value this is a refutation of the instrument. It should not be
taken at face value, and stating why is more useful than the number.

A liquidated wallet no longer holds the position. A registry snapshot taken \emph{after} a
liquidation therefore shows nothing that could have predicted it, and the current snapshot cadence
does not reliably precede the events being scored. The measurement consequently cannot distinguish
\emph{``different populations''} from \emph{``we looked too late''}. Reporting the zero without the
confound would be the more striking result and the less honest one.

\subsection{The experiment that discriminates}

A scheduled snapshot-then-observe cycle---sweep the registry, then scan for liquidations, appending
an overlap reading to the ledger each pass---converts a single confounded reading into a series, and
the series separates the hypotheses. Because sweeps are expensive (Section~\ref{sec:capture}), the
cycle persists last-run times so that a restart does not re-trigger a sweep and a job never overlaps
itself; the cadence is roughly six-hourly snapshots against hourly scans.

\begin{enumerate}
  \item \textbf{Mapped notional fraction climbs off zero.} Per-position scoring works; the map gate's
        predictive criterion becomes evaluable as specified, and the instrument stands.
  \item \textbf{Mapped fraction pinned at zero, notional fraction holding near $22\%$.} The registry
        sees the right wallets but never \emph{before} they are liquidated. Validation must move from
        per-position to cluster-level scoring---did liquidations land in the price buckets the map
        marked, regardless of whose positions they were. A weaker claim, but answerable with the
        registry as it stands.
  \item \textbf{Notional fraction collapses as well.} Neither; the registry needs a third seed drawn
        by \emph{liquidation risk}---margin ratio, or distance to liquidation price, both computable
        from data already held---which targets the population that actually fires, at the cost of
        another sweep.
\end{enumerate}

Only the first outcome validates the instrument as specified. None of the three is obtainable from
further analysis: the discriminating variable is calendar time under capture, which is why capture
runs while the rest is built and why day one of capture is the project's real start date. A
point-in-time series cannot be recovered later.

% ============================================================
\section{Relation to the Companion Papers}
\label{sec:companions}
% ============================================================

This paper is the measurement layer beneath a set of related results, and the division is worth
stating explicitly.

[10] supplies the \emph{model}: the density field, the eight scale-invariant geometric
features, the closed-form domino condition, and the absorption damper. It named the
sample-to-population problem as its binding obstacle and proposed an open-interest rescaling.
Sections~\ref{sec:map}--\ref{sec:mismatch} of the present paper are the first measurements against
that obstacle, and they revise it: the sampling rate is now known, and the rescaling is now known to
be insufficient, for a reason the earlier paper did not anticipate.

[11] establishes, on a different venue and feature set, that a strong and universal predictor
can be economically inert once one conditions on the states in which a passive fill actually
occurs---an adverse-selection barrier that bounds monetization independently of signal quality. That
result is why the present instrument treats execution as a first-class layer with its own simulator
and cost hash rather than as a post-processing step.

[13] develops the analytical contract---the \emph{process} as a unit that certifies whether
a feature carries information about a specified target---and the orthogonality mathematics for
counting a signal book's effective degrees of freedom. The gate ladder of Table~\ref{tab:gates} is
the same discipline applied one level lower, to the data rather than to the signals.

[12] documents the predictive structure of the order book itself, including the finding that
conditioning on directionally-consistent maker fills collapses the information coefficient while
conditioning on the presence of any fill raises it. Together with [11] it establishes the
execution boundary this instrument's allocation and execution layers are designed around.

% ============================================================
\section{Limitations and Open Verification Items}
\label{sec:limits}
% ============================================================

We list what is unverified with the same prominence as what is measured, because the design treats an
unverified constant as a live defect rather than a footnote.

\paragraph{Load-bearing constants not yet confirmed.} The side convention $s$ in
\eqref{eq:coverage}---a factor of two on the headline number. The liquidation formula's residual: the
derivation reproduces the published value for $40.8\%$ of positions, with size-tiered maintenance
margin the prime suspect for the rest. The precise construction of the mark price and which price
triggers liquidation. The funding formula, cadence, and clamp. The rate-limit weight table, which
carries a verification stamp that goes stale at ninety days. The streaming connection limit per IP,
which bounds any future per-user subscription work. Node-data format, retention, and access cost,
which sizes the historical-backfill milestone and therefore the persistence gate.

\paragraph{Known measurement gaps.} The order-book stream's cadence is sparser than its documented
hint implies---$45$ records over $70$ seconds across three coins---so any depth-dependent feature
must wait on a proper cadence measurement. Clock skew flipped sign between capture runs ($-196$\,ms
then $+395$\,ms); both are within the two-second tolerance, but a $\sim$$600$\,ms shift over minutes
is unexplained, and the candidates (NTP discipline versus the venue's timestamp semantics) have
different implications for Definition~\ref{def:causal}.

\paragraph{Structural risks.} Reconstruction may drift from truth between reconciliations; drift
beyond tolerance fails the feed gate, which converts the risk into a detection problem rather than
removing it. Cascade episodes overlap heavily, so any downstream model is exposed to the usual
overlapping-label pathologies and requires purging, embargoing, and uniqueness weighting
[7], with a mandatory naive baseline per model. Crowding is sharper here than
elsewhere for the reason that motivates the whole design: an exactly observable map is exactly
observable to hunters too. Finally, the venue's short history is handled with wider embargoes and
smaller models, never by padding with another venue's data---the first rule again.

% ============================================================
\section{Conclusion}
\label{sec:conclusion}
% ============================================================

We have described an instrument for studying liquidation cascades on a venue where the liquidation
map is observed rather than estimated, and reported what it measured. Two design commitments do the
work. \emph{Exact or nothing} converts unobserved positions into a reported coverage deficit instead
of an imputation, which is what makes the deficit measurable at all. \emph{Gates before models}
converts validation into a refusal that binds every downstream command, which is what makes a failing
verdict survive contact with the desire for a result.

The measurements are these. The feed gate passes twenty integrity and causality checks. Registry
coverage is $24.7$--$38.0\%$ of venue open interest by symbol under the conservative side convention.
The derived liquidation price reproduces the venue's published value to within $10^{-4}$ for $40.8\%$
of $2{,}926$ positions, with an unexplained residual we report rather than model away. And the map
gate fails, because its predictive criterion cannot yet be evaluated and an unevaluable criterion is
not permitted to pass.

The finding we regard as general is the population mismatch. A registry seeded to capture notional
captured $21.9\%$ of liquidated notional but only $2.9\%$ of liquidated wallets, and the gap between
those numbers is invisible to a coverage statistic denominated in notional. Sampling optimized for
mass is mis-specified for events, and no shape-preserving rescaling repairs the difference, because
the sampling is biased in shape rather than merely thinned. We further report that the mapped
fraction is currently zero, and that this figure is confounded by snapshot ordering rather than
informative---a distinction that a scheduled snapshot-then-observe series will resolve in one of three
specified directions, only one of which leaves the instrument's claims intact. We state the
discriminating experiment in advance, in the ledger, so that whichever way it resolves, the record
will show that it was specified before it was run.

% ============================================================
\section*{References}
\addcontentsline{toc}{section}{References}
% ============================================================

\begin{enumerate}
    \item Ackerer, D., Hugonnier, J.\ \& Jermann, U.\ (2025). Perpetual
          futures pricing. \textit{Mathematical Finance}, forthcoming.

    \item Brunnermeier, M.\ K.\ \& Pedersen, L.\ H.\ (2009). Market liquidity and funding
          liquidity. \textit{Review of Financial Studies}, 22(6), 2201--2238.

    \item Cont, R., Kukanov, A.\ \& Stoikov, S.\ (2014). The price impact of order book
          events. \textit{Journal of Financial Econometrics}, 12(1), 47--88.

    \item Easley, D., L\'opez de Prado, M.\ \& O'Hara, M.\ (2012). Flow toxicity and
          liquidity in a high-frequency world. \textit{Review of Financial Studies}, 25(5),
          1457--1493.

    \item Haber, S.\ \& Stornetta, W.\ S.\ (1991). How to time-stamp a
          digital document. \textit{Journal of Cryptology}, 3(2), 99--111.

    \item Kyle, A.\ S.\ (1985). Continuous auctions and insider trading.
          \textit{Econometrica}, 53(6), 1315--1335.

    \item L\'opez de Prado, M.\ (2018). \textit{Advances in Financial
          Machine Learning}. Wiley.

    \item Merkle, R.\ C.\ (1980). Protocols for public key cryptosystems.
          In \textit{Proceedings of the IEEE Symposium on Security and Privacy}, 122--134.

    \item Shleifer, A.\ \& Vishny, R.\ (2011). Fire sales in finance and macroeconomics.
          \textit{Journal of Economic Perspectives}, 25(1), 29--48.

    \item Onat, Y.\ (2026). From liquidation geometry to price movement:
          a spatial-density heatmap and cascade-probability model for cryptocurrency perpetual
          futures. Working paper.

    \item Onat, Y.\ (2026). Autonomous microstructure alpha discovery under
          false-discovery control: order-book imbalance, conditional information coefficients, and
          the adverse-selection barrier in Hyperliquid perpetual futures. Working paper.

    \item Onat, Y.\ (2026). Predictive structure of the perpetual-futures
          limit order book: universality, orthogonality, and the execution boundary. Working paper.

    \item Onat, Y.\ (2026). The process: a first-class analytical unit for
          orthogonal information extraction from market microstructure. Working paper.
\end{enumerate}

\end{document}
